\chaptertitle{Two-Part Invention}

                         Two-Part Invention
                                       or,
                        What the Tortoise Said to Achilles
                                by Lewis Carroll'

Achilles had overtaken the Tortoise, and had seated himself comfortably on its back.
   "So you've got to the end of our race-course?" said the Tortoise. "Even though it
DOES consist of an infinite series of distances? I thought some wiseacre or other had
proved that the thing couldn't be done?"
   "It CAN be done," said Achilles. "It HAS been done! Solvitur ambulando. You see the
distances were constantly DIMINISHING; and so-"
   "But if they had been constantly INCREASING?" the Tortoise interrupted. "How
then?"
   "Then I shouldn't be here," Achilles modestly replied; "and You would have got
several times round the world, by this time!"
   "You flatter me-FLATTEN, I mean," said the Tortoise; "for you ARE a heavy weight,
and NO mistake! Well now, would you like to hear of a race-course, that most people
fancy they can get to the end of in two or three steps, while it REALLY consists of an
infinite number of distances, each one longer than the previous one?"
   "Very much indeed!" said the Grecian warrior, as he drew from his helmet (few
Grecian warriors possessed POCKETS in those days) an enormous note-book and pencil.
"Proceed! And speak SLOWLY, please! SHORTHAND isn't invented yet!"
   "That beautiful First Proposition by Euclid!" the Tortoise murmured dreamily. "You
admire Euclid?"
   "Passionately! So far, at least, as one CAN admire a treatise that won't be published
for some centuries to come!"
   "Well, now, let's take a little bit of the argument in that First Proposition just TWO
steps, and the conclusion drawn from them. Kindly enter them in your note-book. And in
order to refer to them conveniently, let's call them A, B, and Z:

     (A) Things that are equal to the same are equal to each other.
     (B) The two sides of this Triangle are things that are equal to the same.
     (Z) The two sides of this Triangle are equal to each other.

  Readers of Euclid will grant, I suppose, that Z follows logically from A and B, so that
any one who accepts A and B as true, MUST accept Z as true?"
  "Undoubtedly! The youngest child in a High School-as soon as High




Two-Part Invention                                                                    51

Schools are invented, which will not be till some two thousand years later-will grant
THAT."
   "And if some reader had NOT yet accepted A and B as true, he might still accept the
SEQUENCE as a VALID one, I suppose?"
   "No doubt such a reader might exist. He might say, `I accept as true the Hypothetical
Proposition that, IF A and B be true, Z must be true; but I DON'T accept A and B as true.'
Such a reader would do wisely in abandoning Euclid, and taking to football."
   "And might there not ALSO be some reader who would say `I accept A and B as true,
but I DON'T accept the Hypothetical'?"
   "Certainly there might. HE, also, had better take to football."
   "And NEITHER of these readers," the Tortoise continued, "is AS YET under any
logical necessity to accept Z as true?"
   "Quite so," Achilles assented.
   "Well, now, I want you to consider ME as a reader of the SECOND kind, and to force
me, logically, to accept Z as true."
   "A tortoise playing football would be-" Achilles was beginning.
   `-an anomaly, of course," the Tortoise hastily interrupted. "Don't wander from the
point. Let's have Z first, and football afterwards!"
   "I'm to force you to accept Z, am I?" Achilles said musingly. "And your present
position is that you accept A and B, but you DON'T accept the Hypothetical-"
   "Let's call it C," said the Tortoise.
   "-but you DON'T accept

  (C) If A and B are true, Z must be true."

  "That is my present position," said the Tortoise.
  "Then I must ask you to accept C."
  "I'll do so," said the Tortoise, "as soon as you've entered it in that notebook of yours.
What else have you got in it?"
  "Only a few memoranda," said Achilles, nervously fluttering the leaves: "a few
memoranda of-of the battles in which I have distinguished myself!"
  "Plenty of blank leaves, I see!" the Tortoise cheerily remarked. "We shall need them
ALL!" (Achilles shuddered.) "Now write as I dictate:

  (A) Things that are equal to the same are equal to each other.
  (B) The two sides of this Triangle are things that are equal to the same.
  (C) If A and B are true, Z must be true.
  (Z) The two sides of this Triangle are equal to each other."

   "You should call it D, not Z," said Achilles. "It comes NEXT to the other three. If you
accept A and B and C, you MUST accept Z.




Two-Part Invention                                                                      52

   “And why must I?”
   "Because it follows LOGICALLY from them. If A and B and C are true, Z MUST be
true. You can't dispute THAT, I imagine?"
   "If A and B and C are true, Z MUST be true," the Tortoise thoughtfully repeated.
"That's ANOTHER Hypothetical, isn't it? And, if I failed to see its truth, I might accept A
and B and C, and STILL not accept Z, mightn't I?"
   "You might," the candid hero admitted; "though such obtuseness would certainly be
phenomenal. Still, the event is POSSIBLE. So I must ask you to grant ONE more
Hypothetical."
   "Very good, I'm quite willing to grant it, as soon as you've written it down. We will
call it

   (D) If A and B and C are true, Z must be true.

Have you entered that in your note-book?"
    "I HAVE!" Achilles joyfully exclaimed, as he ran the pencil into its sheath. "And at
last we've got to the end of this ideal race-course! Now that you accept A and B and C
and D, OF COURSE you accept Z."
    "Do I?" said the Tortoise innocently. "Let's make that quite clear. I accept A and B and
C and D. Suppose I STILL refused to accept Z?"
    "Then Logic would take you by the throat, and FORCE you to do it!" Achilles
triumphantly replied. "Logic would tell you, `You can't help yourself. Now that you've
accepted A and B and C and D, you MUST accept Z!' So you've no choice, you see.",
    "Whatever LOGIC is good enough to tell me is worth WRITING DOWN," said the
Tortoise. "So enter it in your book, please. We will call it

       (E) If A and B and C and D are true, Z must be true.

   Until I've granted THAT, of course I needn't grant Z. So it's quite a NECESSARY
step, you see?"
   "I see," said Achilles; and there was a touch of sadness in his tone.
   Here the narrator, having pressing business at the Bank, was obliged to leave the
happy pair, and did not again pass the spot until some months afterwards. When he did
so, Achilles was still seated on the back of the much-enduring Tortoise, and was writing
in his notebook, which appeared to be nearly full. The Tortoise was saying, "Have you
got that last step written down? Unless I've lost count, that makes a thousand and one.
There are several millions more to come. And WOULD you mind, as a personal favour,
considering what a lot of instruction this colloquy of ours will provide for the Logicians
of the Nineteenth Century-WOULD you mind adopting a pun that my cousin the Mock-
Turtle will then make, and allowing yourself to be renamed TAUGHT-US?"
   "As you please," replied the weary warrior, in the hollow tones of despair, as he buried
his face in his hands. "Provided that YOU, for YOUR part, will adopt a pun the Mock-
Turtle never made, and allow yourself to be re-named A KILL-EASE!"




Two-Part Invention                                                                       53



                            Meaning and Form
                             in Mathematics.

THIS Two-Part Invention was the inspiration for my two characters. Just as Lewis
Carroll took liberties with Zeno's Tortoise and Achilles, so have I taken liberties with
Lewis Carroll's Tortoise and Achilles. In Carroll's dialogue, the same events take place
over and over again, only each time on a higher and higher level; it is a wonderful
analogue to Bach's Ever-Rising Canon. The Carrollian Dialogue, with its wit subtracted
out, still leaves a deep philosophical problem: Do words and thoughts follow formal
rules, or do they not? That problem is the problem of this book.
In this Chapter and the next, we will look at several new formal systems. This will give
us a much wider perspective on the concept of formal system. By the end of these two
Chapters, you should have quite a good idea of the power of formal systems, and why
they are of interest to mathematicians and logicians.

                                    The pq-System
The formal system of this Chapter is called the pq-system. It is not important to
mathematicians or logicians-in fact, it is just a simple invention of mine. Its importance
lies only in the fact that it provides an excellent example of many ideas that play a large
role in this book. There are three distinct symbols of the pq-system:

                                       p       q    -
-the letters p, q, and the hyphen.
The pq-system has an infinite number of axioms. Since we can't write them all down, we
have to have some other way of describing what they are. Actually, we want more than
just a description of the axioms; we want a way to tell whether some given string is an
axiom or not. A mere description of axioms might characterize them fully and yet
weakly-which was the problem with the way theorems in the MIU-system were
characterized. We don't want to have to struggle for an indeterminate-possibly infinite
length of time, just to find out if some string is an axiom or not. Therefore, we will define
axioms in such a way that there is an obvious decision procedure for axiomhood of a
string composed of p's, q's, and hyphens.




Meaning and Form in Mathematics                                                           54

DEFINITION: xp-qx is an axiom, whenever x is composed of hyphens only.

Note that 'x' must stand for the same string of hyphens in both occurrences For example, -
-p-q---is an axiom. The literal expression `xp-qx-' i,, not an axiom, of course (because `x'
does not belong to the pq-system); it is more like a mold in which all axioms are cast-and
it is called an axiom schema.
         The pq-system has only one rule of production:

RULE: Suppose x, y, and z all stand for particular strings containing only hyphens. And
suppose that x py qz is known to be a theorem. The` xpy-qz- is a theorem.

For example, take x to be'--', y to be'---', and z to be'-'. The rule tells us:

        If --p---q- turns out to be a theorem, then so will --p----q--.


As is typical of rules of production, the statement establishes a causal connection between
the theoremhood of two strings, but without asserting theoremhood for either one on its
own.
        A most useful exercise for you is to find a decision procedure for the theorems of
the pq-system. It is not hard; if you play around for a while you will probably pick it up.
Try it.


                                 The Decision Procedure
I presume you have tried it. First of all, though it may seem too obvious to mention, I
would like to point out that every theorem of the pq-system has three separate groups of
hyphens, and the separating elements are one p, and one q, in that order. (This can be
shown by an argument based on "heredity", just the way one could show that all MIU-
system theorems had to begin with M.) This means that we can rule out, from its form
alone, o string such as --p--p--p--q .
Now, stressing the phrase "from its form alone" may seem silly; what else is there to a
string except its form? What else could possibly play a roll in determining its properties?
Clearly nothing could. But bear this in mint as the discussion of formal systems goes on;
the notion of "form" will star to get rather more complicated and abstract, and we will
have to think more about the meaning of the word "form". In any case, let us give the
name well formed string to any string which begins with a hyphen-group, then ha one p,
then has a second hyphen-group, then a q, and then a final hyphen-group.
Back to the decision procedure ... The criterion for theoremhood is that the first two
hyphen-groups should add up, in length, to the third




Meaning and Form in Mathematics                                                          55

hyphen-group. for instance, --p--q - is a theorem, since 2 plus 2 equals 4, whereas --p--q-
is not, since 2 plus 2 is not 1. To see why this is the proper criterion, look first at the
axiom schema. Obviously, it only manufactures axioms which satisfy the addition
criterion. Second, look at the rule of production. If the first string satisfies the addition
criterion, so must the second one-and conversely, if the first string does not satisfy the
addition criterion, then neither does the second string. The rule makes the addition
criterion into a hereditary property of theorems: any theorem passes the property on to its
offspring. This shows why the addition criterion is correct.
        There is, incidentally, a fact about the pq-system which would enable us to say
with confidence that it has a decision procedure, even before finding the addition
criterion. That fact is that the pq-system is not complicated by the opposing currents of
lengthening and shortening rules; it has only lengthening rules. Any formal system which
tells you how to make longer theorems from shorter ones, but never the reverse, has got
to have a decision procedure for its theorems. For suppose you are given a string. First
check whether it's an axiom or not (I am assuming that there is a decision procedure for
axiomhood-otherwise, things are hopeless). If it is an axiom, then it is by definition a
theorem, and the test is over. So suppose instead that it's not an axiom. Then, to be a
theorem, it must have come from a shorter string, via one of the rules. By going over the
various rules one by one, you can pinpoint not only the rules that could conceivably
produce that string, but also exactly which shorter strings could be its forebears on the
"family tree". In this way, you "reduce" the problem to determining whether any of
several new but shorter strings is a theorem. Each of them can in turn be subjected to the
same test. The worst that can happen is a proliferation of more and more, but shorter and
shorter, strings to test. As you continue inching your way backwards in this fashion, you
must be getting closer to the source of all theorems-the axiom schemata. You just can't
get shorter and shorter indefinitely; therefore, eventually either you will find that one of
your short strings is an axiom, or you'll come to a point where you're stuck, in that none
of your short strings is an axiom, and none of them can be further shortened by running
some rule or other backwards. This points out that there really is not much deep interest
in formal systems with lengthening rules only; it is the interplay of lengthening and
shortening rules that gives formal systems a certain fascination..

                              Bottom-up vs. Top-down

The method above might be called a top-down decision procedure, to be contrasted with a
bottom-up decision procedure, which I give now. It is very reminiscent of the genie's
systematic theorem-generating method for the MIU-system, but is complicated by the
presence of an axiom schema. We are going to form a "bucket" into which we throw
theorems as they are generated. Here is how it is done:




Meaning and Form in Mathematics                                                           56

(1a) Throw the simplest possible axiom (-p-q--) into the bucket.
(I b) Apply the rule of inference to the item in the bucket, and put the result into the
   bucket.
(2a) Throw the second-simplest axiom into the bucket.
(2b) Apply the rule to each item in the bucket, and throw all results into the bucket.
(3a) Throw the third-simplest axiom into the bucket.
(3b) Apply the rule to each item in the bucket, and throw all results into the bucket.

   etc., etc.

A moment's reflection will show that you can't fail to produce every theorem of the pq-
system this way. Moreover, the bucket is getting filled with longer and longer theorems,
as time goes on. It is again a consequence of that lack of shortening rules. So if you have
a particular string, such as --p---q---- , which you want to test for theoremhood, just
follow the numbered steps, checking all the while for the string in question. If it turns up-
theorem! If at some point everything that goes into the bucket is longer than the string in
question, forget it-it is not a theorem. This decision procedure is bottom=up because it is
working its way up from the basics, which is to say the axioms. The previous decision
procedure is top-down because it does precisely the reverse: it works its way back down
towards the basics.

                          Isomorphisms Induce Meaning
Now we come to a central issue of this Chapter-indeed of the book. Perhaps you have
already thought to yourself that the pq-theorems are like additions. The string --p---q--- is
a theorem because 2 plus 3 equals 5. It could even occur to you that the theorem --p---q--
is a statement, written in an odd notation, whose meaning is that 2 plus 3 is 5. Is this a
reasonable way to look at things? Well, I deliberately chose 'p' to remind you of 'plus',
and 'q' to remind you of 'equals' . . . So, does the string --p---q---- actually mean "2 plus
3 equals 5"?
         What would make us feel that way? My answer would be that we have perceived
an isomorphism between pq-theorems and additions. In the Introduction, the word
"isomorphism" was defined as an information preserving transformation. We can now go
into that notion a little more deeply, and see it from another perspective. The word
"isomorphism' applies when two complex structures can be mapped onto each other, in
such a way that to each part of one structure there is a corresponding part in the other
structure, where "corresponding" means that the two part play similar roles in their
respective structures. This usage of the word "isomorphism" is derived from a more
precise notion in mathematics.




Meaning and Form in Mathematics                                                           57

        It is cause for joy when a mathematician discovers an isomorphism between two
structures which he knows. It is often a "bolt from the blue", and a source of
wonderment. The perception of an isomorphism between two known structures is a
significant advance in knowledge-and I claim that it is such perceptions of isomorphism
which create meanings in the minds of people. A final word on the perception of
isomorphisms: since they come in many shapes and sizes, figuratively speaking, it is not
always totally clear when you really have found an isomorphism. Thus, "isomorphism" is
a word with all the usual vagueness of words-which is a defect but an advantage as well.
        In this case, we have an excellent prototype for the concept of isomorphism.
There is a "lower level" of our isomorphism-that is, a mapping between the parts of the
two structures:

                                         p <= => plus
                                        q <= => equals
                                         - <= => one
                                         -- <= => two
                                        ---- <= => three
                                               etc.

This symbol-word correspondence has a name: interpretation.
        Secondly, on a higher level, there is the correspondence between true statements
and theorems. But-note carefully-this higher-level correspondence could not be perceived
without the prior choice of an interpretation for the symbols. Thus it would be more
accurate to describe it as a correspondence between true statements and interpreted
theorems. In any case we have displayed a two-tiered correspondence, which is typical of
all isomorphisms.
        When you confront a formal system you know nothing of, and if you hope to
discover some hidden meaning in it, your problem is how to assign interpretations to its
symbols in a meaningful way-that is, in such a way that a higher-level correspondence
emerges between true statements and theorems. You may make several tentative stabs in
the dark before finding a good set of words to associate with the symbols. It is very
similar to attempts to crack a code, or to decipher inscriptions in an unknown language
like Linear B of Crete: the only way to proceed is by trial and error, based on educated
guesses. When you hit a good choice, a "meaningful" choice, all of a sudden things just
feel right, and work speeds up enormously. Pretty soon everything falls into place. The
excitement of such an experience is captured in The Decipherment of Linear B by John
Chadwick.
        But it is uncommon, to say the least, for someone to be in the position of
"decoding" a formal system turned up in the excavations of a ruined civilization!
Mathematicians (and more recently, linguists, philosophers, and some others) are the only
users of formal systems, and they invariably have an interpretation in mind for the formal
systems which they use and publish. The idea of these people is to set up a formal system
whose




Meaning and Form in Mathematics                                                        58

Theorems reflect some portion of reality isomorphically. In such a case, the choice of
symbols is a highly motivated one, as is the choice of typographical rules of production.
When I devised the pq-system, I was in position. You see why I chose the symbols I
chose. It is no accident theorems are isomorphic to additions; it happened because I
deliberately sought out a way to reflect additions typographically.

                Meaningless and Meaningful Interpretations
        You can choose interpretations other than the one I chose. You need make every
theorem come out true. But there would be very little reason make an interpretation in
which, say, all theorems came out false, certainly even less reason to make an
interpretation under which there is no correlation at all, positive or negative, between
theoremhood and tri Let us therefore make a distinction between two types of
interpretations a formal system. First, we can have a meaningless interpretation, one un
which we fail to see any isomorphic connection between theorems of system, and reality.
Such interpretations abound-any random choice a will do. For instance, take this one:

                                        p <= => horse
                                        q <= => happy
                                        - <= => apple


         Now -p-q-- acquires a new interpretation: "apple horse apple hat apple apple"-a
poetic sentiment, which might appeal to horses, and mi! even lead them to favor this
mode of interpreting pq-strings! However, t interpretation has very little
"meaningfulness"; under interpretative, theorems don't sound any truer, or any better,
than nontheorems. A ho might enjoy "happy happy happy apple horse" (mapped onto q q
q) just as much as any interpreted theorem.
         The other kind of interpretation will be called meaningful. Under si an
interpretation, theorems and truths correspond-that is, an isomorphism exists between
theorems and some portion of reality. That is why it is good to distinguish between
interpretations and meanings. Any old word can be used as an interpretation for `p', but
`plus' is the only meaningful choice we've come up with. In summary, the meaning of `p'
seems to be 'plus’ though it can have a million different interpretations.

                           Active vs. Passive Meanings
Probably the most significant fact of this Chapter, if understood deeply this: the pq-
system seems to force us into recognizing that symbols of a formal system, though
initially without meaning, cannot avoid taking on "meaning" of sorts at least if an
isomorphism is found. The difference between meaning it formal system and in a
language is a very important one, however. It is this:




Meaning and Form in Mathematics                                                       59

in a language, when we have learned a meaning for a word, we then mar-c new
statements based on the meaning of the word. In a sense the meaning becomes active,
since it brings into being a new rule for creating sentences. This means that our command
of language is not like a finished product: the rules for making sentences increase when
we learn new meanings. On the other hand, in a formal system, the theorems are
predefined, by the rules of production. We can choose "meanings" based on an
isomorphism (if we can find one) between theorems and true statements. But this does
not give us the license to go out and add new theorems to the established theorems. That
is what the Requirement of Formality in Chapter I was warning you of.
        In the MIU-system, of course, there was no temptation to go beyond the four
rules, because no interpretation was sought or found. But here, in our new system, one
might be seduced by the newly found "meaning" of each symbol into thinking that the
string

                                        --p--p--p--q

is a theorem. At least, one might wish that this string were a theorem. But wishing doesn't
change the fact that it isn't. And it would be a serious mistake to think that it "must" be a
theorem, just because 2 plus 2 plus 2 plus 2 equals 8. It would even be misleading to
attribute it any meaning at all, since it is not well-formed, and our meaningful
interpretation is entirely derived from looking at well-formed strings.
        In a formal system, the meaning must remain passive; we can read each string
according to the meanings of its constituent symbols, but we do not have the right to
create new theorems purely on the basis of the meanings we've assigned the symbols.
Interpreted formal systems straddle the line between systems without meaning, and
systems with meaning. Their strings can be thought of as "expressing" things, but this
must come only as a consequence of the formal properties of the system.

                                  Double-Entendre!
And now, I want to destroy any illusion about having found the meanings for the symbols
of the pq-system. Consider the following association:

                                      p <= => equals
                                      q <= => taken from
                                      - <= => one
                                      -- <= => two
                                            etc.

Now, --p---q---- has a new interpretation: "2 equals 3 taken from 5". Of course it is a true
statement. All theorems will come out true under this new interpretation. It is just as
meaningful as the old one. Obviously, it is silly to ask, "But which one is the meaning of
the string?" An interpreta




Meaning and Form in Mathematics                                                           60

tion will me meaningful to the extent that it accurately reflects some isomorphism to the
real world. When different aspects of the real world a isomorphic to each other (in this
case, additions and subtractions), or single formal system can be isomorphic to both, and
therefore can take ( two passive meanings. This kind of double-valuedness of symbols at
strings is an extremely important phenomenon. Here it seems trivial curious, annoying.
But it will come back in deeper contexts and bring with it a great richness of ideas.
        Here is a summary of our observations about the pq-system. Und either of the two
meaningful interpretations given, every well-form( string has a grammatical assertion for
its counterpart-some are true, son false. The idea of well formed strings in any formal
system is that they a those strings which, when interpreted symbol for symbol, yield
grammatical sentences. (Of course, it depends on the interpretation, but usually, there one
in mind.) Among the well-formed strings occur the theorems. The: are defined by an
axiom schema, and a rule of production. My goal in inventing the pq-system was to
imitate additions: I wanted every theorem] to express a true addition under interpretation;
conversely, I wanted every true addition of precisely two positive integers to be
translatable into a string, which would be a theorem. That goal was achieved. Notice,
then fore, that all false additions, such as "2 plus 3 equals 6", are mapped into strings
which are well-formed, but which are not theorems.


                            Formal Systems and Reality
This is our first example of 'a case where a formal system is based upon portion of
reality, and seems to mimic it perfectly, in that its theorems a] isomorphic to truths about
that part of reality. However, reality and tt formal system are independent. Nobody need
be aware that there is a isomorphism between the two. Each side stands by itself-one plus
or equals two, whether or not we know that -p-q-- is a theorem; and -p-q-- is still a
theorem whether or not we connect it with addition.
         You might wonder whether making this formal system, or any form system, sheds
new light on truths in the domain of its interpretation. Hat we learned any new additions
by producing pq-theorems? Certainly not but we have learned something about the nature
of addition as process-namely, that it is easily mimicked by a typographical rule
governing meaningless symbols. This still should not be a big surprise sing addition is
such a simple concept. It is a commonplace that addition can I captured in the spinning
gears of a device like a cash register.
         But it is clear that we have hardly scratched the surface, as far formal systems go;
it is natural to wonder about what portion of reality co be imitated in its behavior by a set
of meaningless symbols governed I formal rules. Can all of reality be turned into a formal
system? In a very broad sense, the answer might appear to be yes. One could suggest, for
instance, that reality is itself nothing but one very complicated formal




Meaning and Form in Mathematics                                                           61

system. Its symbols do not move around on paper, but rather in a three-dimensional
vacuum (space); they are the elementary particles of which everything is composed.
(Tacit assumption: that there is an end to the descending chain of matter, so that the
expression "elementary particles" makes sense.) The "typographical rules" are the laws of
physics, which tell how, given the positions and velocities of all particles at a given
instant, to modify them, resulting in a new set of positions and velocities belonging to the
"next" instant. So the theorems of this grand formal system are the possible
configurations of particles at different times in the history of the universe. The sole axiom
is (or perhaps, was) the original configuration of all the particles at the "beginning of
time". This is so grandiose a conception, however, that it has only the most theoretical
interest; and besides, quantum mechanics (and other parts of physics) casts at least some
doubt on even the theoretical worth of this idea. Basically, we are asking if the universe
operates deterministically, which is an open question.

                        Mathematics and Symbol Manipulation
         Instead of dealing with such a big picture, let's limit ourselves to mathematics as
our "real world". Here, a serious question arises: How can we be sure, if we've tried to
model a formal system on some part of mathematics, that we've done the job accurately-
especially if we're not one hundred per cent familiar with that portion of mathematics
already? Suppose the goal of the formal system is to bring us new knowledge in that
discipline. How will we know that the interpretation of every theorem is true, unless
we've proven that the isomorphism is perfect? And how will we prove that the
isomorphism is perfect, if we don't already know all about the truths in the discipline to
begin with?
         Suppose that in an excavation somewhere, we actually did discover some
mysterious formal system. We would try out various interpretations and perhaps
eventually hit upon one which seemed to make every theorem come out true, and every
nontheorem come out false. But this is something which we could only check directly in
a finite number of cases. The set of theorems is most likely infinite. How will we know
that all theorems express truths under this interpretation, unless we know everything there
is to know about both the formal system and the corresponding domain of interpretation?
         It is in somewhat this odd position that we will find ourselves when we attempt to
match the reality of natural numbers (i.e., the nonnegative integers: 0, 1, 2, ...) with the
typographical symbols of a formal system. We will try to understand the relationship
between what we call "truth" in number theory and what we can get at by symbol
manipulation.
         So let us briefly look at the basis for calling some statements of number theory
true, and others false. How much is 12 times 12? Everyone knows it is 144. But how
many of the people who give that answer have actually at




Meaning and Form in Mathematics                                                           62

any time in their lives drawn a 12 by 12 rectangle, and then counted the little squares in
it? Most people would regard the drawing and counting unnecessary. They would instead
offer as proof a few marks on paper, such as are shown below:

                                                X 12
                                                ------
                                                ------


        And that would be the "proof". Nearly everyone believes that if you counted the
squares, you would get 144 of them; few people feel that outcome is in doubt.
        The conflict between the two points of view comes into sharper focus when you
consider the problem of determining the value 987654321 x 123456789. First of all, it is
virtually impossible to construct the appropriate rectangle; and what is worse, even if it
were constructed and huge armies of people spent centuries counting the little squares, o
a very gullible person would be willing to believe their final answer. It is just too likely
that somewhere, somehow, somebody bobbled just a little bit. So is it ever possible to
know what the answer is? If you trust the symbolic process which involves manipulating
digits according to certain simple rules, yes. That process is presented to children as a
device which gets right answer; lost in the shuffle, for many children, are the rhyme
reason of that process. The digit-shunting laws for multiplication are based mostly on a
few properties of addition and multiplication which are assumed to hold for all numbers.

                               The Basic Laws of Arithmetic
The kind of assumption I mean is illustrated below. Suppose that you down a few sticks:

                                    / // // // / /

Now you count them. At the same time, somebody else counts them, starting from the
other end. Is it clear that the two of you will get the s: answer? The result of a counting
process is independent of the way in which it is done. This is really an assumption about
what counting i would be senseless to try to prove it, because it is so basic; either you s or
you don't-but in the latter case, a proof won't help you a bit.
        From this kind of assumption, one can get to the commutativity and associativity
of addition (i.e., first that b + c = c + b always, and second that b + (c + d) = (b + c) + d
always). The same assumption can also you to the commutativity and associativity of
multiplication; just think of




Meaning and Form in Mathematics                                                            63

many cubes assembled to form a large rectangular solid. Multiplicative commutativity
and associativity are just the assumptions that when you rotate the solid in various ways,
the number of cubes will not change. Now these assumptions are not verifiable in all
possible cases, because the number of such cases is infinite. We take them for granted;
we believe them (if we ever think about them) as deeply as we could believe anything.
The amount of money in our pocket will not change as we walk down the street, jostling
it up and down; the number of books we have will not change if we pack them up in a
box, load them into our car, drive one hundred miles, unload the box, unpack it, and place
the books in a new shelf. All of this is part of what we mean by number.
        There are certain types of people who, as soon as some undeniable fact is written
down, find it amusing to show why that "fact" is false after all. I am such a person, and as
soon as I had written down the examples above involving sticks, money, and books, I
invented situations in which they were wrong. You may have done the same. It goes to
show that numbers as abstractions are really quite different from the everyday numbers
which we use.
        People enjoy inventing slogans which violate basic arithmetic but which illustrate
"deeper" truths, such as "1 and 1 make 1" (for lovers), or "1 plus 1 plus 1 equals 1" (the
Trinity). You can easily pick holes in those slogans, showing why, for instance, using the
plus-sign is inappropriate in both cases. But such cases proliferate. Two raindrops
running down a windowpane merge; does one plus one make one? A cloud breaks up into
two clouds-more evidence for the same? It is not at all easy to draw a sharp line between
cases where what is happening could be called "addition", and where some other word is
wanted. If you think about the question, you will probably come up with some criterion
involving separation of the objects in space, and making sure each one is clearly
distinguishable from all the others. But then how could one count ideas? Or the number
of gases comprising the atmosphere? Somewhere, if you try to look it up, you can
probably find a statement such as, "There are 17 languages in India, and 462 dialects."
There is something strange about precise statements like that, when the concepts
"language" and "dialect" are themselves fuzzy.

                                        Ideal Numbers
        Numbers as realities misbehave. However, there is an ancient and innate sense in
people that numbers ought not to misbehave. There is something clean and pure in the
abstract notion of number, removed from counting beads, dialects, or clouds; and there
ought to be a way of talking about numbers without always having the silliness of reality
come in and intrude. The hard-edged rules that govern "ideal" numbers constitute
arithmetic, and their more advanced consequences constitute number theory. There is
only one relevant question to be asked, in making the transition from numbers as
practical things to numbers as formal things. Once you have




Meaning and Form in Mathematics                                                          64

                FIGURE 13. Liberation, by M.C. Escher (lithograph, 1955).




Meaning and Form in Mathematics                                             65

decided to try to capsulize all of number theory in an ideal system, is it really possible to
do the job completely? Are numbers so clean and crystalline and regular that their nature
can be completely captured in the rules of a formal system? The picture Liberation (Fig.
13), one of Escher's most beautiful, is a marvelous contrast between the formal and the
informal, with a fascinating transition region. Are numbers really as free as birds? Do
they suffer as much from being crystallized into a rule-obeying system? Is there a
magical transition region between numbers in reality and numbers on paper?
        When I speak of the properties of natural numbers, I don't just mean properties
such as the sum of a particular pair of integers. That can be found out by counting, and
anybody who has grown up in this century cannot doubt the mechanizability of such
processes as counting, adding, multiplying, and so on. I mean the kinds of properties
which mathematicians are interested in exploring, questions for which no counting-
process is sufficient to provide the answer-not even theoretically sufficient. Let us take a
classic example of such a property of natural numbers. The statement is: "There are
infinitely many prime numbers." First of all, there is no counting process which will ever
be able to confirm, or refute, this assertion. The best we could do would be to count
primes for a while and concede that there are "a lot". But no amount of counting alone
would ever resolve the question of whether the number of primes is finite or infinite.
There could always be more. The statement-and it is called "Euclid's Theorem" (notice
the capital "T")-is quite unobvious. It may seem reasonable, or appealing, but it is not
obvious. However, mathematicians since Euclid have always called it true. What is the
reason?

                                         Euclid's Proof
The reason is that reasoning tells them it is so. Let us follow the reasoning involved. We
will look at a variant of Euclid's proof. This proof works by showing that whatever
number you pick, there is a prime larger than it. Pick a number-N. Multiply all the
positive integers starting with 1 and ending with N; in other words, form the factorial of
N, written "N!". What you get is divisible by every number up to N. When you add 1 to
N!, the result

            can't be a multiple of 2 (because it leaves 1 over, when you divide
                                     by 2);
            can't be a multiple of 3 (because it leaves I over, when you divide
                                     by 3);
            can't be a multiple of 4 (because it leaves 1 over, when you divide
                                     by 4);
            .
            .
            .




Meaning and Form in Mathematics                                                           66

            can't be a multiple of N (because it leaves 1 over, when you
                                  divide by N);

        In other words, N! + 1, if it is divisible at all (other than by 1 and itself only is
divisible by numbers greater than N. So either it is itself prime, or prime divisors are
greater than N. But in either case we've shown the must exist a prime above N. The
process holds no matter what number is. Whatever N is, there is a prime greater than N.
And thus ends the demonstration of the infinitude of the primes.
        This last step, incidentally, is called generalization, and we will meet again later
in a more formal context. It is where we phrase an argument terms of a single number
(N), and then point out that N was unspecified and therefore the argument is a general
one.
        Euclid's proof is typical of what constitutes "real mathematics". It simple,
compelling, and beautiful. It illustrates that by taking several rash short steps one can get
a long way from one's starting point. In our case, t starting points are basic ideas about
multiplication and division and forth. The short steps are the steps of reasoning. And
though eve individual step of the reasoning seems obvious, the end result is not obvious.
We can never check directly whether the statement is true or not; } we believe it, because
we believe in reasoning. If you accept reasoning there seems to be no escape route; once
you agree to hear Euclid out, you’ll have to agree with his conclusion. That's most
fortunate-because it mea that mathematicians will always agree on what statements to
label "true and what statements to label "false".
        This proof exemplifies an orderly thought process. Each statement related to
previous ones in an irresistible way. This is why it is called "proof'' rather than just "good
evidence". In mathematics the goal always to give an ironclad proof for some unobvious
statement. The very fact of the steps being linked together in an ironclad way suggests ti
there may be a patterned structure binding these statements together. TI structure can
best be exposed by finding a new vocabulary-a stylized vocabulary, consisting of
symbols-suitable only for expressing statements about numbers. Then we can look at the
proof as it exists in its translated version. It will be a set of statements which are related,
line by line, in some detectable way. But the statements, since they're represented by
means a small and stylized set of symbols, take on the aspect of patterns. In other words,
though when read aloud, they seem to be statements about numb and their properties, still
when looked at on paper, they seem to be abstract patterns-and the line-by-line structure
of the proof may start to look like slow transformation of patterns according to some few
typographical rules.

                                   Getting Around Infinity
Although Euclid's proof is a proof that all numbers have a certain property it avoids
treating each of the infinitely many cases separately. It gets around




Meaning and Form in Mathematics                                                             67

it by using phrases like "whatever N is", or "no matter what number N is". We could also
phrase-the proof over again, so that it uses the phrase "all N". By knowing the appropriate
context and correct ways of using such phrases, we never have to deal with infinitely
many statements. We deal with just two or three concepts, such as the word "all"-which,
though themselves finite, embody an infinitude; and by using them, we sidestep the
apparent problem that there are an infinite number of facts we want to prove.
         We use the word "all" in a few ways which are defined by the thought processes
of reasoning. That is, there are rules which our usage of "all" obeys. We may be
unconscious of them, and tend to claim we operate on the basis of the meaning of the
word; but that, after all, is only a circumlocution for saying that we are guided by rules
which we never make explicit. We have used words all our lives in certain patterns, and
instead of calling the patterns "rules", we attribute the courses of our thought processes to
the "meanings" of words. That discovery was a crucial recognition in the long path
towards the formalization of number theory.
         If we were to delve into Euclid's proof more and more carefully, we would see
that it is composed of many, many small-almost infinitesimal steps. If all those steps were
written out line after line, the proof would appear incredibly complicated. To our minds it
is clearest when several steps are telescoped together, to form one single sentence. If we
tried to look at the proof in slow motion, we would begin to discern individual frames. In
other words, the dissection can go only so far, and then we hit the "atomic" nature of
reasoning processes. A proof can be broken down into a series of tiny but discontinuous
jumps which seem to flow smoothly when perceived from a higher vantage point. In
Chapter VIII, I will show one way of breaking the proof into atomic units, and you will
see how incredibly many steps are involved. Perhaps it should not surprise you, though.
The operations in Euclid's brain when he invented the proof must have involved millions
of neurons (nerve cells), many of which fired several hundred times in a single second.
The mere utterance of a sentence involves hundreds of thousands of neurons. If Euclid's
thoughts were that complicated, it makes sense for his proof to contain a huge number of
steps! (There may be little direct connection between the neural actions in his brain, and a
proof in our formal system, but the complexities of the two are comparable. It is as if
nature wants the complexity of the proof of the infinitude of primes to be conserved, even
when the systems involved are very different from each other.)
         In Chapters to come, we will lay out a formal system that (1) includes a stylized
vocabulary in which all statements about natural numbers can be expressed, and (2) has
rules corresponding to all the types of reasoning which seem necessary. A very important
question will be whether the rules for symbol manipulation which we have then
formulated are really of equal power (as far as number theory is concerned) to our usual
mental reasoning abilities-or, more generally, whether it is theoretically possible to attain
the level of our thinking abilities, by using some formal system.




Meaning and Form in Mathematics                                                           68

                             Sonata
                   for Unaccompanied Achilles

                    The telephone rings; Achilles picks it up.

Achilles: Hello, this is Achilles.
Achilles: Oh, hello, Mr. T. How are you?
Achilles: A torticollis? Oh, I'm sorry to hear it. Do you have any idea what caused it?
Achilles: How long did you hold it in that position?
Achilles: Well, no wonder it's stiff, then. What on earth induced you keep your neck
       twisted that way for so long?
Achilles: Wondrous many of them, eh? What kinds, for example? Achilles: What do you
       mean, "phantasmagorical beasts"?

              FIGURE 14. Mosaic II, by M. C. Escher (lithograph, 1957).




Sonata for Unaccompanied Achilles                                                   69

Achilles: Wasn't it terrifying to see so many of them at the same time? Achilles: A
       guitar!? Of all things to be in the midst of all those weird creatures. Say, don't you
       play the guitar?
Achilles: Oh, well, it's all the same to me.
Achilles: You're right; I wonder why I never noticed that difference between fiddles and
       guitars before. Speaking of fiddling, how would you like to come over and listen
       to one of the sonatas for unaccompanied violin by your favorite composer, J. S.
       Bach? I just bought a marvelous recording of them. I still can't get over the way
       Bach uses a single violin to create a piece with such interest.
Achilles: A headache too? That's a shame. Perhaps you should just go to bed.
Achilles: I see. Have you tried counting sheep?
Achilles: Oh, oh, I see. Yes, I fully know what you mean. Well, if it's THAT distracting,
       perhaps you'd better tell it to me, and let me try to work on it, too.
Achilles: A word with the letters `A', `D', `A', `C' consecutively inside it ... Hmm ...
       What about "abracadabra"?
Achilles: True, "ADAC" occurs backwards, not forwards, in that word. Achilles: Hours
       and hours? It sounds like I'm in for a long puzzle, then. Where did you hear this
       infernal riddle?
Achilles: You mean he looked like he was meditating on esoteric Buddhist matters, but in
       reality he was just trying to think up complex word puzzles?
Achilles: Aha!-the snail knew what this fellow was up to. But how did you come to talk
       to the snail?
Achilles: Say, I once heard a word puzzle a little bit like this one. Do you want to hear it?
       Or would it just drive you further into distraction? Achilles: I agree-can't do any
       harm. Here it is: What's a word that begins with the letters "HE" and also ends
       with "HE"?
Achilles: Very ingenious-but that's almost cheating. It's certainly not what I meant!
Achilles: Of course you're right-it fulfills the conditions, but it's a sort of "degenerate"
       solution. There's another solution which I had in mind. Achilles: That's exactly it!
       How did you come up with it so fast? Achilles: So here's a case where having a
       headache actually might have helped you, rather than hindering you. Excellent!
       But I'm still in the dark on your "ADAC" puzzle.
Achilles: Congratulations! Now maybe you'll be able to get to sleep! So tell me, what is
       the solution?
Achilles: Well, normally I don't like hints, but all right. What's your hint? Achilles: I
       don't know what you mean by "figure" and "ground" in this case.
Achilles: Certainly I know Mosaic II! I know ALL of Escher's works. After all, he's my
       favorite artist. In any case, I've got a print of Mosaic II hanging on my wall, in
       plain view from here.




Sonata for Unaccompanied Achilles                                                         70

Achilles:: Yes, t see all the black animals.
Achilles: Yes, I also see how their "negative space" -- what's left out-- defines the white
       animals.
Achilles: So THAT'S what you mean by "figure" and "ground". But what does that have
       to do with the "ADAC" puzzle?
Achilles: Oh, this is too tricky for me. I think I'M starting to get a headache
Achilles: You want to come over now? But I thought--
Achilles: Very well. Perhaps by then I'll have thought of the right answer to YOUR
       puzzle, using your figure-ground hint, relating it to MY puzzle
Achilles: I'd love to play them for you.
Achilles: You've invented a theory about them?
Achilles: Accompanied by what instrument?
Achilles: Well, if that's the case, it seems a little strange that he would have written out
       the harpsichord part, then, and had it published a s well.
Achilles: I see -- sort of an optional feature. One could listen to them either way -- with
       or without accompaniment. But how would one know what the accompaniment is
       supposed to sound like?
Achilles: Ah, yes, I guess that it is best, after all, to leave it to the listener’s imagination.
       And perhaps, as you said, Bach never even had accompaniment in mind at all.
       Those sonatas seem to work very indeed as they are.
Achilles: Right. Well, I'll see you shortly.
Achilles: Good-bye, Mr. T.




Sonata for Unaccompanied Achilles                                                             71


                          Figure and Ground
                                Primes vs. Composites
THERE IS A strangeness to the idea that concepts can be captured by simple
typographical manipulations. The one concept so far captured is that of addition, and it
may not have appeared very strange. But suppose the goal were to create a formal system
with theorems of the form Px, the letter `x' standing for a hyphen-string, and where the
only such theorems would be ones in which the hyphen-string contained exactly a prime
number of hyphens. Thus, P--- would be a theorem, but P---- would not. How could this
be done typographically? First, it is important to specify clearly what is meant by
typographical operations. The complete repertoire has been presented in the MIU-system
and the pq-system, so we really only need to make a list of the kinds of things we have
permitted:

      (1) reading and recognizing any of a finite set of symbols;
      (2) writing down any symbol belonging to that set;
      (3) copying any of those symbols from one place to another;
      (4) erasing any of those symbols;
      (5) checking to see whether one symbol is the same as another;
      (6) keeping and using a list of previously generated theorems.

The list is a little redundant, but no matter. What is important is that it clearly involves
only trivial abilities, each of them far less than the ability to distinguish primes from
nonprimes. How, then, could we compound some of these operations to make a formal
system in which primes are distinguished from composite numbers?

                                     The tq-System
A first step might be to try to solve a simpler, but related, problem. We could try to make
a system similar to the pq-system, except that it represents multiplication, instead of
addition. Let's call it the tq-system, `t' for times'. More specifically, suppose X, Y, and Z
are, respectively, the numbers of hyphens in the hyphen-strings x, y, and z. (Notice I am
taking special pains to distinguish between a string and the number of hyphens it
contains.) Then we wish the string x ty q z to be a theorem if and only if X times Y
equals Z. For instance, --t---q----- should be a theorem because 2 times 3 equals 6, but --
t--q--- should not be a theorem. The tq-system can be characterized just about as easily as
the pq-system namely, by using just one axiom schema and one rule of inference:




Figure and Ground                                                                         72

AXIOM SCHEMA: xt-qx is an axiom, whenever x is a hyphen string.
.
RULE OF INFERENCE: Suppose that x, y, and z are all hyphen-strings. An suppose that
     x ty qz is an old theorem. Then, xty-qzx is a ne' theorem.

Below is the derivation of the theorem --t---q-----
:
     (1)       --t-q--         (axiom)
     (2)       --t—q----       (by rule of inference, using line (1) as the old theorem)
     (3)       --t---q ------- (by rule of inference, using line (2) as the old theorem)

Notice how the middle hyphen-string grows by one hyphen each time the rule of
inference is applied; so it is predictable that if you want a theorem with ten hyphens in
the middle, you apply the rule of inference nine times in a row.

                             Capturing Compositeness
Multiplication, a slightly trickier concept than addition, has now bee] "captured"
typographically, like the birds in Escher's Liberation. What about primeness? Here's a
plan that might seem smart: using the tq-system define a new set of theorems of the form
Cx, which characterize compost. numbers, as follows:

RULE: Suppose x, y, and z are hyphen-strings. If x-ty-qz is a theorem then C z is a
theorem.

This works by saying that Z (the number of hyphens in z) is composite a long as it is the
product of two numbers greater than 1-namely, X + (the number of hyphens in x-), and Y
+ 1 (the number of hyphens in y I am defending this new rule by giving you some
"Intelligent mode justifications for it. That is because you are a human being, and want t,
know why there is such a rule. If you were operating exclusively in the "Mechanical
mode", you would not need any justification, since M-mod. workers just follow the rules
mechanically and happily, never questioning; them!
        Because you work in the I-mode, you will tend to blur in your mind the
distinction between strings and their interpretations. You see, things Cal become quite
confusing as soon as you perceive "meaning" in the symbol which you are manipulating.
You have to fight your own self to keep from thinking that the string'---' is the number 3.
The Requirement of Formality, which in Chapter I probably seemed puzzling (because it
seemed so obvious), here becomes tricky, and crucial. It is the essential thing which
keeps you from mixing up the I-mode with the M-mode; or said another way, it keeps
you from mixing up arithmetical facts with typographical theorems.




Figure and Ground                                                                          73

                          Illegally Characterizing Primes
It is very tempting to jump from the C-type theorems directly to P-type theorems, by
proposing a rule of the following kind:

PROPOSED RULE: Suppose x is a hyphen-string. If Cx is not a theorem, then Px is a
theorem.
The fatal flaw here is that checking whether Cx is not a theorem is not an explicitly
typographical operation. To know for sure that MU is not a theorem of the MIU-system,
you have to go outside of the system ... and so it is with this Proposed Rule. It is a rule
which violates the whole idea of formal systems, in that it asks you to operate informally-
that is, outside the system. Typographical operation (6) allows you to look into the
stockpile of previously found theorems, but this Proposed Rule is asking you to look into
a hypothetical "Table of Nontheorems". But in order to generate such a table, you would
have to do some reasoning outside the system-reasoning which shows why various strings
cannot be generated inside the system. Now it may well be that there is another formal
system which can generate the "Table of Nontheorems", by purely typographical means.
In fact, our aim is to find just such a system. But the Proposed Rule is not a typographical
rule, and must be dropped.
         This is such an important point that we might dwell on it a bit more. In our C-
system (which includes the tq-system and the rule which defines C-type theorems), we
have theorems of the form Cx, with `x' standing, as usual, for a hyphen-string. There are
also nontheorems of the form Cx. (These are what I mean when I refer to "nontheorems",
although of course tt-Cqq and other ill-formed messes are also nontheorems.) The
difference is that theorems have a composite number of hyphens, nontheorems have a
prime number of hyphens. Now the theorems all have a common "form", that is, originate
from a common set of typographical rules. Do all nontheorems also have a common
"form", in the same sense? Below is a list of C-type theorems, shown without their
derivations. The parenthesized numbers following them simply count the hyphens in
them.

                         C---- (4)
                         C-------- (6)
                         C----------------(8)
                         C-----------------(9)
                         C--------------------(10)
                         C -------------------- (12)
                         C------------------------ (14)
                         C ------------------------ (15)
                         C----------------------------(16)
                         C----------------------------(18)


                                           .
                                           .
                                           .

Figure and Ground                                                                        74

I he "holes" in this list are the nontheorems. I o repeat the earlier quest Do the holes also
have some "form" in common? Would it be reasonable say that merely by virtue of being
the holes in this list, they share a common form? Yes and no. That they share some
typographical quality is and able, but whether we want to call it "form" is unclear. The
reason hesitating is that the holes are only negatively defined-they are the things that are
left out of a list which is positively defined.

       Figure and Ground

This recalls the famous artistic distinction between figure and ground. When a figure or
"positive space" (e.g., a human form, or a letter, or a still life is drawn inside a frame, an
unavoidable consequence is that its complementary shape-also called the "ground", or
"background", or "negative space"-has also been drawn. In most drawings, however, this
fig ground relationship plays little role. The artist is much less interested in ground than
in the figure. But sometimes, an artist will take interest in ground as well.
         There are beautiful alphabets which play with this figure-ground distinction. A
message written in such an alphabet is shown below. At fir looks like a collection of
somewhat random blobs, but if you step back ways and stare at it for a while, all of a
sudden, you will see seven letters appear in this ..
         .




                                            FIGURE 15.


For a similar effect, take a look at my drawing Smoke Signal (Fig. 139). Along these
lines, you might consider this puzzle: can you somehow create a drawing containing
words in both the figure and the ground?
        Let us now officially distinguish between two kinds of figures: cursively
drawable ones, and recursive ones (by the way, these are my own terms are not in
common usage). A cursively drawable figure is one whose ground is merely an
accidental by-product of the drawing act. A recursive figure is one whose ground can be
seen as a figure in its own right. Usually this is quite deliberate on the part of the artist.
The "re" in "recursive" represents the fact that both foreground and background are
cursively drawable – the figure is "twice-cursive". Each figure-ground boundary in a
recursive figure is a double-edged sword. M. C. Escher was a master at drawing recursive
figures-see, for instance, his beautiful recursive drawing of birds (Fig. 16).




Figure and Ground                                                                          75

FIGURE 16. Tiling of the plane using birds, by M. C. Escher (from a 1942 notebook).

        Our distinction is not as rigorous as one in mathematics, for who can definitively
say that a particular ground is not a figure? Once pointed out, almost any ground has
interest of its own. In that sense, every figure is recursive. But that is not what I intended
by the term. There is a natural and intuitive notion of recognizable forms. Are both the
foreground and background recognizable forms? If so, then the drawing is recursive. If
you look at the grounds of most line drawings, you will find them rather unrecognizable.
This demonstrates that

      There exist recognizable forms whose negative space is not any recognizable form.

In more "technical" terminology, this becomes:

       There exist cursively drawable figures which are not recursive.

       Scott Kim's solution to the above puzzle, which I call his "FIGURE-FIGURE
Figure", is shown in Figure 17. If you read both black and white,




Figure and Ground                                                                          76

              FIGURE 17. FIGURE-FIGURE Figure, by Scott E. Kim (1975).




Figure and Ground                                                        77

you will see "FIGURE" everywhere, but "GROUND" nowhere! It is a paragon of
recursive figures. In this clever drawing, there are two nonequivalent ways of
characterizing the black regions:

  (1) as the negative space to the white regions;
  (2) as altered copies of the white regions (produced by coloring and shifting each
        white region).

(In the special case of the FIGURE-FIGURE Figure, the two characterizations are
equivalent-but in most black-and-white pictures, they would not be.) Now in Chapter
VIII, when we create our Typographical Number Theory (TNT), it will be our hope that
the set of all false statements of number theory can be characterized in two analogous
ways:

     (1) as the negative space to the set of all TNT-theorems;
     (2) as altered copies of the set of all TNT-theorems (produced by negating each
           TNT-theorem).

       But this hope will be dashed, because:

        (1) inside the set of all nontheorems are found some truths
        (2) outside the set of all negated theorems are found some falsehoods
        .
You will see why and how this happens, in Chapter XIV. Meanwhile, ponder over a
pictorial representation of the situation (Fig. 18).

                              Figure and Ground in Music
One may also look for figures and grounds in music. One analogue is the distinction
between melody and accompaniment-for the melody is always in the forefront of our
attention, and the accompaniment is subsidiary, in some sense. Therefore it is surprising
when we find, in the lower lines of a piece of music, recognizable melodies. This does
not happen too often in post-baroque music. Usually the harmonies are not thought of as
foreground. But in baroque music-in Bach above all-the distinct lines, whether high or
low or in between, all act as "figures". In this sense, pieces by Bach can be called
"recursive".
        Another figure-ground distinction exists in music: that between on-beat and off-
beat. If you count notes in a measure "one-and, two-and, three-and, four-and", most
melody-notes will come on numbers, not on "and"'s. But sometimes, a melody will be
deliberately pushed onto the "and" 's, for the sheer effect of it. This occurs in several
etudes for the piano by Chopin, for instance. It also occurs in Bach-particularly in his
Sonatas and Partitas for unaccompanied violin, and his Suites for unaccompanied cello.
There, Bach manages to get two or more musical lines going simultaneously. Sometimes
he does this by having the solo instrument play "double stops"-two notes at once. Other
times, however, he




Figure and Ground                                                                     78

FIGURE 18. Considerable visual symbolism is featured in this diagram of the relation
between various classes of TNT strings. The biggest box represents the set of all TNT
strings The next-biggest box represents the set of all well-formed TNT strings. Within it is
found~ set of all sentences of TNT. Now things begin to get interesting. The set of
theorems pictured as a tree growing out of a trunk (representing the set of axioms). The
tree-symbol chosen because of the recursive growth pattern which it exhibits: new
branches (theorems constantly sprouting from old ones. The fingerlike branches probe
into the corners of constraining region (the set of truths), yet can never fully occupy it.
The boundary beta the set of truths and the set of falsities is meant to suggest a randomly
meandering coastline which, no matter how closely you examine it, always has finer
levels of structure, an consequently impossible to describe exactly in any finite way. (See
B. Mandelbrot's book Fractals.) The reflected tree represents the set of negations of
theorems: all of them false yet unable collectively to span the space of false statements.
[Drawing by the author.]

puts one voice on the on-beats, and the other voice on the off-beats, so ear separates them
and hears two distinct melodies weaving in and out, - harmonizing with each other.
Needless to say, Bach didn't stop at this level of complexity...

                    Recursively Enumerable Sets vs. Recursive Sets
        Now let us carry back the notions of figure and ground to the domain formal
systems. In our example, the role of positive space is played by C-type theorems, and the
role of negative space is played by strings with a




Figure and Ground                                                                        79

prime number of hyphens. So far, the only way we have found to represent prime
numbers typographically is as a negative space. Is there, however, some way-I don't care
how complicated-of representing the primes as a positive space-that is, as a set of
theorems of some formal system?
       Different people's intuitions give different answers here. I remember quite vividly
how puzzled and intrigued I was upon realizing the difference between a positive
characterization and a negative characterization. I was quite convinced that not only the
primes, but any set of numbers which could be represented negatively, could also be
represented positively. The intuition underlying my belief is represented by the question:
"How could a figure and its ground not carry exactly the same information?" They
seemed to me to embody the same information, just coded in two complementary ways.
What seems right to you?
       It turns out I was right about the primes, but wrong in general. This astonished
me, and continues to astonish me even today. It is a fact that:

      There exist formal systems whose negative space (set of nontheorems) is not
      the positive space (set of theorems) of any formal system.

        This result, it turns out, is of depth equal to Gödel’s Theorem-so it is not
surprising that my intuition was upset. I, just like the mathematicians of the early
twentieth century, expected the world of formal systems and natural numbers to be more
predictable than it is. In more technical terminology, this becomes:

       There exist recursively enumerable sets which are not recursive.

The phrase recursively enumerable (often abbreviated "r.e.") is the mathematical
counterpart to our artistic notion of "cursively drawable"-and recursive is the counterpart
of "recursive". For a set of strings to be "r.e." means that it can be generated according to
typographical rules-for example, the set of C-type theorems, the set of theorems of the
MIU-system-indeed, the set of theorems of any formal system. This could be compared
with the conception of a "figure" as "a set of lines which can be generated according to
artistic rules" (whatever that might mean!). And a "recursive set" is like a figure whose
ground is also a figure-not only is it r.e., but its complement is also r.e.
         It follows from the above result that:

      There exist formal systems for which there is no typographical decision
      procedure.

       How does this follow? Very simply. A typographical decision procedure is a
method which tells theorems from nontheorems. The existence of such a test allows us to
generate all nontheorems systematically, simply by going down a list of all strings and
performing the test on them one at a time, discarding ill-formed strings and theorems
along the way. This amounts to




Figure and Ground                                                                         80

a typographical method for generating the set of nontheorems. But according to the
earlier statement (which we here accept on faith), for some systems this is not possible.
So we must conclude that typographical decision procedures do not exist for all formal
systems.
        Suppose we found a set F of natural numbers (`F' for `Figure') whi4 we could
generate in some formal way-like the composite numbers. Suppose its complement is the
set G (for 'Ground')-like the primes. Together F and G make up all the natural numbers,
and we know a rule for making all the numbers in set F, but we know no such rule for
making all tl numbers in set G. It is important to understand that if the members of were
always generated in order of increasing size, then we could always characterize G. The
problem is that many r.e. sets are generated I methods which throw in elements in an
arbitrary order, so you never know if a number which has been skipped over for a long
time will get included you just wait a little longer.
        We answered no to the artistic question, "Are all figures recursive We have now
seen that we must likewise answer no to the analogous question in mathematics: "Are all
sets recursive?" With this perspective, 1 us now come back to the elusive word "form".
Let us take our figure-set and our ground-set G again. We can agree that all the numbers
in set have some common "form"-but can the same be said about numbers in s G? It is a
strange question. When we are dealing with an infinite set to sta with-the natural
numbers-the holes created by removing some subs may be very hard to define in any
explicit way. And so it may be that th< are not connected by any common attribute or
"form". In the last analysis it is a matter of taste whether you want to use the word
"form"-but just thinking about it is provocative. Perhaps it is best not to define "form", bi
to leave it with some intuitive fluidity.
        Here is a puzzle to think about in connection with the above matter Can you
characterize the following set of integers (or its negative space)

       1       3       7      12      18      26      35     45      56      69...

How is this sequence like the FIGURE-FIGURE Figure?


                           Primes as Figure Rather than Ground
Finally, what about a formal system for generating primes? How is it don< The trick is to
skip right over multiplication, and to go directly to nondivisibility as the thing to
represent positively. Here are an axiom schema and rule for producing theorems which
represent the notion that one number does not divide (D N D) another number exactly:

       AXIOM SCHEMA: xy D N Dx where x and y are hyphen-strings.

       For example ----D N D--, where x has been replaced by'--'and y by ‘---“.




Figure and Ground                                                                         81

RULE: If x D N Dy is a theorem, then so is x D N Dx y.

If you use the rule twice, you can generate this theorem:

               -----D N D --------------

which is interpreted as "5 does not divide 12". But ---D N D------------ is not a theorem.
What goes wrong if you try to produce it?
         Now in order to determine that a given number is prime, we have to build up
some knowledge about its nondivisibility properties. In particular, we want to know that
it is not divisible by 2 or 3 or 4, etc., all the way up to 1 less than the number itself. But
we can't be so vague in formal systems as to say "et cetera". We must spell things out.
We would like to have a way of saying, in the language of the system, "the number Z is
divisor free up to X", meaning that no number between 2 and X divides Z. This can be
done, but there is a trick to it. Think about it if you want.
         Here is the solution:

RULE: If --D N D z is a theorem, so is z D F--.

RULE: If z D Fx is a theorem and also x-D N Dz is a theorem, z D Fx- is a theorem.

These two rules capture the notion of divisor freeness. All we need to do is to say that
primes are numbers which are divisor-free up to 1 less than themselves:

RULE: If z-DFz is a theorem, then Pz- is a theorem.

Oh-let's not forget that 2 is prime!

Axiom: P--.

And there you have it. The principle of representing primality formally is that there is a
test for divisibility which can be done without any backtracking. You march steadily
upward, testing first for divisibility by 2, then by 3, and so on. It is this "monotonicity" or
unidirectionality-this absence of cross-play between lengthening and shortening,
increasing and decreasing-that allows primality to be captured. And it is this potential
complexity of formal systems to involve arbitrary amounts of backwards-forwards
interference that is responsible for such limitative results as Gödel’s Theorem, Turing's
Halting Problem, and the fact that not all recursively enumerable sets are recursive.




Figure and Ground                                                                           82

